\newpage
\phantomsection
\label{prf:polynomials-are-continuous}

\begin{remark*}[Return]
\hyperref[thm:polynomials-are-continuous]{Return to Theorem}
\end{remark*}

\begin{theorem*}[Any Real Polynomial Is a Continuous Function]
Let
\[
  P(x) = a_0 + a_1 x + a_2 x^2 + \cdots + a_n x^n
\]
be a real polynomial. Then $P$ is a continuous function.
\end{theorem*}

\begin{proof}
\textbf{Professional Standard Proof.}~\\
Write $P(x) = T_0 + T_1 + \cdots + T_n$ where $T_k = a_k x^k$ for $k = 0, 1, \ldots, n$. Each term $T_k$ is continuous: if $k = 0$, then $T_0 = a_0$ is constant; if $k \geq 1$ and $a_k \neq 0$, then $T_k = a_k x^k$ is a scalar multiple of the continuous power function $x^k$. By the constant multiple rule, each $T_k$ is continuous. Since $P$ is a finite sum of continuous functions, $P$ is continuous by the sum rule.
\end{proof}

\begin{proof}
\textbf{Detailed Learning Proof.}~\\

\textbf{Step 1.} Decompose the polynomial into individual terms. Write $P(x) = T_0 + T_1 + \cdots + T_n$ where $T_k = a_k x^k$ for each $k = 0, 1, \ldots, n$. This expresses $P$ as a finite sum of monomial terms, allowing analysis of each component separately.

\textbf{Step 2.} Establish continuity of each term $T_k$. For $k = 0$, the term $T_0 = a_0$ is a constant function, hence continuous. For $k \geq 1$, if $a_k = 0$, then $T_k = 0$ is constant and continuous. If $a_k \neq 0$, then $T_k = a_k x^k$ is a scalar multiple of the power function $x^k$, which is continuous. By the constant multiple rule for continuous functions, $T_k$ is continuous.

\textbf{Step 3.} Apply the sum rule to conclude continuity of $P$. Since each term $T_k$ is continuous and $P$ is the finite sum $T_0 + T_1 + \cdots + T_n$, the sum rule for continuous functions guarantees that $P$ is continuous.
\end{proof}

\begin{remark*}[Proof structure]
This is a direct proof by decomposition. The strategy decomposes the polynomial into monomial terms, establishes continuity of each term using basic continuity properties (constants and power functions), then applies the sum rule. The finite nature of the sum ensures the argument closes cleanly without issues of uniform convergence.
\end{remark*}

\begin{remark*}[Dependencies]~\\
\begin{itemize}
  \item \hyperref[lem:power-function-continuous]{Power Function Is Continuous}
  \item \hyperref[thm:constant-multiple-rule-continuous]{Constant Multiple Rule for Continuous Functions}
  \item \hyperref[thm:sum-rule-continuous-functions]{Sum Rule for Continuous Functions}
\end{itemize}
\end{remark*}
